\chapter{Prorrogação do prazo de entrega do diploma}
\label{chap:pf-prorroga_diploma}

\section{O que é?}

É a extensão em seis meses do prazo de entrega do diploma, para alteração da inscrição de provisória para definitiva.

\section{Quem pode utilizar esse serviço?}

Psicólogas/os que, tendo solicitado o diploma em tempo hábil, tenham o documento retido ou não expedido pela entidade formadora devido a
débitos em aberto.

\section{Como solicitar o serviço?}

Se não estiver em posse do diploma até o vencimento da inscrição provisória, solicite a prorrogação pelo período de até seis meses,
enviando o protocolo ou a declaração do pedido de emissão do diploma fornecido pela entidade formadora (IES) para a sua subsede.

\section{Que informações ou documentos são necessários?}

Protocolo ou declaração do pedido de emissão do diploma fornecido pela entidade formadora (IES)

\section{Quais são as etapas para a realização desse serviço?}

\begin{enumerate}
	\item
	A/o psicóloga/o solicita a prorrogação do prazo de entrega do diploma
	antes do prazo de dois anos, contados da solicitação de inscrição
	provisória;
	\item
	o CRP~SP analisa a documentação enviada e concede a extensão do prazo.
\end{enumerate}

\subsection{Se a/o psicóloga/o apresenta o diploma no prazo de seis meses}

\begin{enumerate}
	\setcounter{enumi}{2}
	\item
	\hyperref[entrega-da-carteira-de-identidade-profissional-cip]{A
		carteira de identidade profissional é entregue}.
\end{enumerate}

\subsection{Se a/o psicóloga/o não apresenta o diploma no prazo de seis meses}

\begin{enumerate}
	\setcounter{enumi}{2}
	\item
	O CRP~SP notifica a/o psicóloga/o do vencimento do prazo concedendo 30 dias para a entrega do documento ou nova e última prorrogação do prazo de entrega do documento;
	\item
	a/o psicóloga/o encaminha nova solicitação de prorrogação, junto com o protocolo de solicitação da entidade formadora;
	\item
	se o diploma for entregue no prazo da segunda prorrogação, a \hyperref[chap:pf-cip]{CIP definitiva é expedida e entregue}. Caso contrário, o CRP~SP procede ao cancelamento da inscrição provisória.
\end{enumerate}

\section{Quanto custa?}

Não tem custo.

\section{Quanto tempo leva?}

Após a submissão do pedido no \emph{site}, o prazo é de até 30 dias.

\section{Como ter acesso ao atual \emph{status} do serviço?}

No \emph{site} do CRP~SP, acessando a página da/o psicóloga/o com seu \emph{login} e senha, no campo \href{https://cfpsp.brctotal.com/crp06_servicosonline/pgsRequerimento/MeusRequerimentos.aspx}{Meus requerimentos}.

\section{Legislação relacionada}

\begin{itemize}
	\item
	\href{https://atosoficiais.com.br/cfp/resolucao-administrativa-financeira-n-3-2007-institui-a-consolidacao-das-resolucoes-do-conselho-federal-de-psicologia?origin=instituicao&q=resolu\%C3\%A7\%C3\%A3o\%203/2007}{Resolução
		CFP nº~03/2007} (Consolidação das Resoluções do Conselho Federal de
	Psicologia);
	\item
	\href{https://transparencia.cfp.org.br/wp-content/uploads/sites/7/2020/05/Manual-de-Procedimentos-Administrativos-e-Financeiro-CFP-RES-20.2018.pdf}{\href{https://atosoficiais.com.br/cfp/resolucao-administrativa-financeira-n-20-2018-altera-o-manual-de-procedimentos-administrativo-e-financeiro-do-sistema-conselhos-de-psicologia-anexo-da-resolucao-cfp-n-202018-a-resolucao-cfp-n-03-2007-a-resolucao-cfp-n-16-2019-e-da-outras-providencias}{Resolução
			CFP nº~20/2018}}, que determina a revisão e ampliação do
	\href{https://transparencia.cfp.org.br/wp-content/uploads/sites/7/2020/05/Manual-de-Procedimentos-Administrativos-e-Financeiro-CFP-RES-20.2018.pdf}{\emph{Manual
			de procedimentos administrativos e financeiros}}.
\end{itemize}

\section{Serviços correlatos}

\begin{itemize}
	\item
	\hyperref[chap:pf-cip]{Entrega da carteira de identidade profissional};
	\item
	\hyperref[chap:pf-inscricao_definitiva]{alteração da inscrição de provisória para definitiva}.
\end{itemize}

\sumario